\chapter{Verification method}\label{ch:method}

A tracking code is verified against references whose error is known, with instruments whose own error is measured.
The benchmarks follow five rules.

\section{Register before running}
Each case states its question, metrics, configuration, references and success criteria in \nolinkurl{benchmarks/DESIGN.md}
before any run.
Criteria fixed after seeing results are labelled as amendments, with the reason.
A criterion that fails stays failed in the record, even when the reason is benign.
Two examples from the record:
\begin{itemize}
\item Case A asked OPAL RK4 to converge at fourth order in the orbit radius.
It converged at fifth order, as \cref{eq:rk4-decay} shows it must.
The criterion was wrong, and it is recorded as failed, with the derivation added after the run.
\item Case H asked that the Hermite surface have value jumps below $10^{-10}$ across cell edges, probed at $\pm10^{-9}$ of a
cell.
A smooth surface changes by its gradient times the probe distance, so the test could not pass as posed.
It was replaced by a scaling test (jumps proportional to the probe distance), documented as a change made after the
result.
\end{itemize}

\section{References with a known error}
Three kinds of reference were used, from strongest to weakest:
\begin{enumerate}
\item \textbf{Closed forms.}
A uniform field (case A): circle, period, and the Boris polygon of \cref{eq:boris-radius}.
A field linear in $r$ (case B0): closed orbit from $P/|q| = |B_0|\,r\,(1-\kappa(r-r_0))$, field index $n = \kappa r/(1-\kappa(r-r_0))$,
tunes from \cref{eq:weak-tunes}.
\item \textbf{Manufactured solutions.}
FROM\_ORBIT builds a field from a prescribed family of isochronous closed orbits \cite{planche2019,planche2023}.
The orbits and the revolution frequency are then exact by construction.
Its tunes come from a geometric formula; case G1 step 1 measured them against RF-Track on a smooth interpolation of the
map (\cref{sec:caseG1}).
\item \textbf{External codes.}
FIXPO closed orbits and tunes for the PSI Ring, OPAL-cycl \cite{adelmann2019} and RF-Track
\cite{latina2023}.
Their error is not known in advance; it is bounded by comparing them with each other.
\end{enumerate}

\section{Field parity}
Codes compared on ``the same field'' must see the same field.
In the synthetic-field cases (A, B0, B2, D2), every converted field is sampled from the LIONS field evaluation itself,
so a conversion cannot add a second field model, and before each run the field of every code is compared with LIONS at
random points.
On the external maps (G1, G2) this is not possible: OPAL reads the PSI file natively and recomputes the derivatives,
and RF-Track runs on splines; there the field models differ by design and the comparison measures that.
In case B0, $B_z$ agrees within \BzeroParityBz~T and $B_x$, $B_y$ within \BzeroParityBxy~T, below the derived bound of
the LIONS chord error.
A failed parity check stops the run.
One did: the first case B2 check, written for \num{1e-13}~T on all components, failed at \num{4.9e-5}~T, and the bound
was rederived before running (\texttt{DESIGN.md}, section 14).

\section{Separate the error sources}
A tracking result mixes three errors: the field model (interpolation), the integrator (step), and the measurement
(section crossing, finite differences, Newton tolerance).
The cases are built to change one at a time:
\begin{itemize}
\item same field, different integrators and steps: B0, D2;
\item same integrator, same nodes, different interpolation: G1 step 3 (bilinear against splines, Boris fixed);
\item same field, different integrator: G2 step 2c (RF-Track on the bilinear field);
\item same code path, two implementations: the numpy Boris harness reproduces LIONS to \GValidationMax{} over one turn before it is used
for anything else.
\end{itemize}

\section{Competing hypotheses}
When a result is unexpected, the rival explanations are written down with the observation that would separate them,
and a discriminating run is registered.
\Cref{tab:evidence} shows the evidence matrix that built up for the bilinear-map noise across cases B, G1 and G2.

\begin{table}[t]
\centering\small
\caption{Evidence matrix for the noise of LIONS one-turn maps on bilinear maps.
H1: the bilinear field (kinks at cell edges). H1b: the Boris step on a kinked field. H2: a LIONS-specific defect.
H3: the finite-difference Jacobian. H4: a near-integer tune amplifying small errors.
$+$ supports, $-$ contradicts, $\cdot$ neutral.}
\label{tab:evidence}
\begin{tabular}{p{7.4cm}ccccc}
\toprule
Observation & H1 & H1b & H2 & H3 & H4\\
\midrule
B0 (bilinear exact in $B_z$): orbit order 2, $\det-1\sim10^{-9}$ & $+$ & $\cdot$ & $-$ & $-$ & $\cdot$\\
B: RF-Track on the closed-form field is clean at the same tune & $+$ & $\cdot$ & $\cdot$ & $\cdot$ & $\cdot$\\
B-gradient: OPAL spread follows $1/(2\sin\pi\Qr)$, LIONS partly & $\cdot$ & $\cdot$ & $\cdot$ & $\cdot$ & $+/-$\\
B: halving the map step does not reduce the scatter & $-$ (as $O(h^2)$) & $\cdot$ & $\cdot$ & $\cdot$ & $\cdot$\\
G1 step 3: same Boris, same nodes, spline field removes the floor & $+$ & $+$ & $\cdot$ & $\cdot$ & $\cdot$\\
G2 2b: numpy Boris on the bilinear field equals LIONS & $+$ & $+$ & $-$ & $\cdot$ & $\cdot$\\
G2 2b: $\det$ changes at most \GTwoBJacVar\,\% from $h=10^{-7}$ to $10^{-5}$ m & $\cdot$ & $\cdot$ & $\cdot$ & $-$ & $\cdot$\\
G2 2c: RF-Track on the bilinear field, $|\det-1|$ \GTwoCRatioMin{} to \GTwoCRatioMax{} times smaller & $\cdot$ & $+$ & $\cdot$ & $\cdot$ & $\cdot$\\
H: F1 (C$^1$) with Boris removes the $\det$ floor on all maps & $+$ & $+$ & $\cdot$ & $\cdot$ & $\cdot$\\
\bottomrule
\end{tabular}
\end{table}

The current best explanation is H1b: the Boris step crossing gradient jumps of a bilinear field.
H1 as first stated, an $O(h^2)$ interpolation error, failed its own prediction.
In case B, halving the azimuthal step left the scatter unchanged (ratio \BHalfDegQrRatioMin{} to
\BHalfDegQrRatioMax), where $O(h^2)$ predicts 0.25.
Kinks of any spacing (H1$'$) remain consistent with every observation.
A direct test, a Jacobian step scan with the closed orbit moved across a cell edge, has not been run.
